\subsection{Perspectives: Towards a Derivation of the Proton-to-Electron Mass Ratio}
  \label{subsec:perspectives_mass_spectrum}

  The proton is modeled, at the exploratory level,
  as a composite bound state within a $Q=3$ sector,
  motivated by generic stability features observed
  in nonlinear solitonic systems~\cite{BattyeSutcliffe2022, MantonSutcliffe2004}.

  At the schematic level,
  \begin{equation}
    \frac{m_p}{m_e}
    \;\sim\;
    \sqrt{\frac{\lambda_{\mathrm{bind}}}{\lambda_e}},
  \end{equation}
  where $\lambda_e$ denotes the elementary-sector stiffness scale
  and $\lambda_{\mathrm{bind}}$ a characteristic composite-sector scale.

  Using the empirical value $m_p/m_e \simeq 1836$
  as an external constraint implies an order-of-magnitude hierarchy
  \begin{equation}
    \frac{\lambda_{\mathrm{bind}}}{\lambda_e}
    \sim 10^{6}.
  \end{equation}
  This relation should be understood solely as a target
  hierarchy for any explicit realization of the mechanism.
  No claim is made that the ratio is predicted
  or uniquely determined at the present stage.

  \subsubsection*{Working Ansatz for $V(\chi)$}
    \label{subsec:explicit_Vchi_ansatz}

    A minimal two-scale effective potential:
    \begin{equation}
      V(\chi)
      = \frac{\lambda}{4}
      (\chi^2 - \eta^2)^2
      + \varepsilon\,\eta^4
      \left[
        1 - \cos\!\left(
                    \frac{\chi}{\eta}
      \right)
      \right],
      \label{eq:Vchi_multiwell_periodic}
    \end{equation}
    with $\varepsilon \ll 1$.

  \subsubsection*{Linking $(\lambda,\eta)$ to Observables}
    \label{subsec:lambda_eta_to_observables}

    Dimensionless control combinations:
    \begin{equation}
      g \equiv \lambda\,\chi_c^2\,\eta^2,
      \quad
      u \equiv \varepsilon.
      \label{eq:dimensionless_controls}
    \end{equation}

    \paragraph{Matching strategy.}
      \label{matching-strategy-using-the-proton-to-electron-ratio}
      The feasibility condition:
    \begin{equation}
      \exists\,(g,u)
      \text{ s.t.}\quad
      \frac{\lambda_{\mathrm{bind}}(Q{=}3;\,g,u)}
      {\lambda_e(g,u)}
      \sim 10^6,
      \label{eq:feasibility_condition_lambda_eta}
      \end{equation}
      while remaining stable under small $(g,u)$ variations.

  \subsubsection*{Numerical Program}
    \label{subsec:numerical_program_mass_ratio}

    The pipeline tests whether the admissibility constraints on~$\Pi$
    together with $V(\chi)$ admit stable sectors with
    $\lambda_{\mathrm{bind}}/\lambda_e \sim 10^6$ without fine-tuning:
    dynamics and formation, soliton harvesting by topological
    classification, stability operator extraction, spectral ratio test,
    and fine-structure diagnostics.
